% =============================================================================
%  CAT1: Characterisation — skeleton paper
%  STATUS: UNPROVEN. Skeleton only; depends on S1 and N1 first.
% =============================================================================
\documentclass[11pt]{article}

\usepackage[utf8]{inputenc}
\usepackage[T1]{fontenc}
\usepackage[expansion=false]{microtype}
\usepackage{amsmath, amssymb, amsthm, mathtools}
\usepackage[margin=1in]{geometry}
\usepackage{enumitem}
\usepackage[hidelinks]{hyperref}

\theoremstyle{plain}
\newtheorem{theorem}{Theorem}
\newtheorem*{namedconjecture}{Conjecture}
\theoremstyle{definition}
\newtheorem{definition}[theorem]{Definition}

\title{Characterisation of $\mathbf{LifeCat}$ via the Decomposition Functor \\
       \large (UNPROVEN --- skeleton paper for CAT1)}
\author{Liam McGuire\thanks{Unaffiliated researcher, Seattle, USA. liammcg44@gmail.com}}
\date{\today}

\begin{document}
\maketitle

\begin{abstract}
This paper is a \textbf{skeleton} that records the statement of the
characterisation conjecture CAT1 from the scaffolding
paper~\cite{emergent_systems_2026}. \textbf{No proof is attempted.}
The result is UNPROVEN. CAT1 depends on S1 and N1, both of which are
also UNPROVEN; CAT1 should be the \emph{last} of the top-level
conjectures attempted, not the first.
\end{abstract}

\section{Statement (UNPROVEN)}

\begin{quote}
\emph{``CAT1 (characterisation): $\mathbf{LifeCat} \simeq
\mathrm{image}(\mathrm{Decomp})$, with the vitality-profile equivalence
preserved.''}
\end{quote}
\noindent (Scaffolding paper~\cite{emergent_systems_2026},
\S\texttt{sec:cat1}.)

\section{Plan of attack (no proof attempted)}

CAT1 is essentially S1 + N1 + a coherence theorem. Once S1 and N1
land, the remaining content is:

\begin{enumerate}[leftmargin=2em,itemsep=4pt]
  \item \textbf{Equivalence direction $\subseteq$ (UNPROVEN).} Every
    pair in $\mathbf{LifeCat}$ is in $\mathrm{image}(\mathrm{Decomp})$.
    Follows from N1 plus a verification that the decomposed system
    remains in $\mathbf{LifeCat}$ (preserves the membership predicate).
  \item \textbf{Equivalence direction $\supseteq$ (UNPROVEN).} Every
    object in $\mathrm{image}(\mathrm{Decomp})$ is in
    $\mathbf{LifeCat}$ when paired with an appropriate observer.
    Follows from S1 applied per object.
  \item \textbf{Vitality-profile preservation (UNPROVEN).} The
    equivalence respects the profile attribute: $(k, \boldsymbol{\sigma})$
    is the same on both sides of the equivalence. This is the
    coherence-theorem step.
\end{enumerate}

\section{Dependencies}

CAT1 should not be attempted before S1 and N1 land. The paper exists
as a placeholder so the framework's characterisation claim is on
record and can be cited as a target.

\section{Status and follow-up}

\textbf{Current status: UNPROVEN. Blocked on S1 and N1.}

\bibliographystyle{plain}
\bibliography{references}

\end{document}
